\section{Del 2: DNBs beslutning og de ansattes motivasjon}

Motivasjon på arbeidsplassen defineres i organisasjonsteorien som de biologiske, psykologiske og sosiale faktorene som aktiverer, gir retning til og opprettholder målrettet atferd. I kunnskapsintensive virksomheter er medarbeidernes drivkraft avgjørende for verdiskaping og trivsel. Når DNB avskaffer faste hjemmekontordager og krever at fravær må avklares med nærmeste leder, griper dette direkte inn i de ansattes arbeidsbetingelser. For å forklare hvorfor tiltaket har utløst det Finansforbundet beskriver som en \enquote{uvanlig sterk respons}, må beslutningen belyses gjennom selvbestemmelsesteori, behovs- og prestasjonsmotiver, rettferdighetsteori samt teorier om psykologiske kontrakter.

\subsection{Selvbestemmelsesteori og jobbengasjement}

Selvbestemmelsesteorien (SDT) til Deci og Ryan postulerer at best mulig motivasjon, funksjon og trivsel avhenger av tre grunnleggende psykologiske behov: autonomi (selvbestemmelse), kompetanse (mestring) og tilhørighet (relasjoner) \parencite[s.~101--102]{nerstad2023}. Autonomi innebærer å kunne bestemme over egen tid, ta egne valg og utøve selvstendig initiativ \parencite[s.~101--102]{nerstad2023}. Hjemmekontor har i praksis fungert som en sentral arena for en slik opplevd autonomi.

Innstrammingen i DNB synliggjør et markant paradoks innenfor SDT: Ledelsen begrunner tiltaket med et ønske om å styrke felles læring (kompetanse) og fellesskap (tilhørighet), men forsøker å oppnå dette ved å innskrenke de ansattes autonomi. Casen viser til solid organisasjonspsykologisk støtte for at autonomi er avgjørende for motivasjon og prestasjon. Når hjemmearbeid underlegges godkjenning, flyttes motivasjonen fra indre, selvbestemt regulering mot kontrollert, ytre regulering. Uten reell valgfrihet reduseres medarbeidernes evne til å internalisere verdien av samvær; det som kunne vært opplevd som fellesskap, oppfattes i stedet som kontroll, noe som svekker arbeidsgleden og fremkaller psykologisk reaktans (motstand).

Dette påvirker også de ansattes jobbengasjement. Gjensidig tillit, variasjon og opplevd autonomi er forutsetninger for et engasjerende arbeidsmiljø \parencite[s.~152]{einarsen2023}. Mens motivasjon representerer den indre psykologiske kraften, utgjør jobbengasjement den observerbare energien og innsatsen som rettes mot arbeidsoppgavene \parencite[s.~152]{einarsen2023}. Ved å erstatte generell tillit med individuell godkjenning signaliseres mistillit. Konsekvensen er at både den underliggende indre motivasjonen og det observerbare engasjementet svekkes \parencite[s.~152]{einarsen2023}.

\subsection{Behovsteorier, tilhørighet og prestasjonsmotiver}

DNBs ledelse argumenterer imidlertid med at økt fysisk tilstedeværelse styrker arbeidsmiljøet og samhandlingen. Dette argumentet har forankring i klassisk motivasjonsteori. I Maslows behovshierarki og Alderfers ERG-teori er sosiale behov og tilhørighet avgjørende for menneskets funksjon \parencite[s.~97--98]{nerstad2023}. Mens Maslow hevder at sosiale tilhørighetsbehov må dekkes før vekstbehov kan realiseres på toppen av hierarkiet, postulerer Alderfer at tilhørighets- og vekstbehov kan opptre samtidig \parencite[s.~97--98]{nerstad2023}. Vekstbehovet omfatter ønsket om å utvikle personlige evner og realisere sitt fulle potensial i arbeidet \parencite[s.~97--98]{nerstad2023}. Fysisk samhandling er en sentral arena for å dekke disse behovene. Nora Dakos i Vend bemerker at mye hjemmekontor for nyutdannede kan føre til ensomhet og tap av uformell læring fra erfarne kolleger \parencite{ingebrigtsen_vend_2026}, mens Silje Marie Anfindsen i DNV opplever dager med tomt kontor som \enquote{litt sosialt kjedelig} \parencite{ingebrigtsen_dnv_2026}. Prestasjonsmotivteori underbygger dette: Vår atferd drives av de tre ubevisste motivene prestasjon, tilhørighet og makt. Disse motivene stimuleres og forsterkes gjennom mellommenneskelig kontakt og synlighet på kontoret \parencite[s.~100]{nerstad2023}. DNBs ønske om å samle medarbeiderne hviler dermed på et teoretisk premiss om at fysisk nærhet fremmer relasjonsbygging og kollektiv kompetanseheving.

\subsection{Rettferdighetsteori, prosedyrerettferdighet og likeverdsspenning}

Misnøyen oppstår primært fordi virkemiddelet oppleves som dypt urettferdig. J.~Stacy Adams' rettferdighetsteori forklarer hvordan arbeidstakere vurderer bytteforholdet mellom egen innsats (input) og hva de mottar av goder og anerkjennelse (outcome) \parencite[s.~105]{nerstad2023}.

Når forholdstallet oppleves som ubalansert, oppstår likeverdsspenning \parencite[s.~105]{nerstad2023}. HR-direktør Gro Gotteberg i DNV uttaler treffende: \enquote{Jeg tror trusselen for oss, er å ta bort noe som er ansett som et gode} \parencite{ingebrigtsen_dnv_2026}. Ved å fjerne faste hjemmekontordager krever DNB økt innsats (input) i form av tapt reisetid og mer krevende hverdagslogistikk (for eksempel levering av barn), samtidig som et etablert gode (outcome) i form av tidsfleksibilitet fjernes.

Som Pål Behrens i Finansforbundet påpeker, forsterkes denne ubalansen av at banken har innført \enquote{free seating} og \enquote{clean desk} \parencite{biernat2026}. De ansatte pålegges økt oppmøteplikt, men returnerer til et kontorlandskap uten en fast eller personlig arbeidsstasjon, noe som ytterligere reduserer det opplevde utbyttet (outcome). For å utjevne likeverdsspenningen kan ansatte ifølge Adams redusere egen innsats, fremme lønnskrav eller si opp \parencite[s.~105]{nerstad2023}. Nora Dakos i Vend bekrefter dette ved å antyde at bortfall av fleksibilitet ville vært en mulig \enquote{dealbreaker} for hennes ansettelsesforhold \parencite{ingebrigtsen_vend_2026}. Videre forutsetter Adams at individer foretar sosiale sammenligninger med eksterne referansegrupper \parencite[s.~105]{nerstad2023}. Når DNB-ansatte observerer at aktører som Vend opprettholder stor frihet, forsterkes opplevelsen av urettferdighet og relativ deprivasjon.

I tillegg rammes opplevelsen av prosedyrerettferdighet, altså hvorvidt beslutningsprosessene oppleves som konsistente, transparente og inkluderende. Ved å overlate avgjørelsen om hjemmekontor til den enkelte linjeleders skjønn, innfører DNB en ordning med stor fare for vilkårlighet og avdelingsvise forskjeller. Dette skaper uforutsigbarhet, svekker legitimiteten til prosessen og svekker tilliten til ledelsessystemet.

\subsection{Psykologiske kontrakter, hygienefaktorer og ledelsesrammer}

Reaksjonene kan også analyseres gjennom Herzbergs tofaktorteori. Mens motivasjonsfaktorer (som anerkjennelse og ansvar) fremmer varig trivsel, representerer hygienefaktorer (som fysiske arbeidsbetingelser og bedriftens retningslinjer) grunnleggende forventninger. Hygienefaktorer skaper i seg selv sjelden entusiasme, men fravær eller forringelse av dem utløser akutt misnøye. Fleksibilitet i form av hjemmekontor har gradvis etablert seg som en hygienefaktor i moderne kunnskapsbedrifter. Når godet innskrenkes, skaper det umiddelbar demotivasjon.

Dette henger tett sammen med den psykologiske kontrakten, de uskrevne, gjensidige forventningene mellom arbeidsgiver og arbeidstaker. Etter pandemien oppstod en implisitt avtale: Ansatte leverte dokumentert høy produktivitet mot at ledelsen innvilget fleksibilitet og autonomi. DNBs ensidige instruks oppfattes som et kontraktsbrudd. Forskning indikerer at slike brudd fører til kynisme, redusert lojalitet, økt turnover og en nedgang i organisatorisk medborgerskapsatferd (OCB), den frivillige ekstrainnsatsen som opprettholder et velfungerende kunnskapsmiljø.

Samtidig må situasjonen nyanseres i tråd med poengene til Myhre og Aas \parencite{myhre2026}: Autonomi betyr ikke grenseløs individuell valgfrihet. Tydelig ledelse, krav og felles rammer er nødvendig for å koordinere komplekse organisasjoner og motvirke kognitive blindsoner som uhensiktsmessig multitasking. Likevel feiler DNB i gjennomføringen: I stedet for å forankre tiltaket som en meningsfull felles standard gjennom medvirkning, fremstår endringen som en kontrollorientert innstramming.

For å gjenreise jobbengasjementet må ledelsen aktivt stimulere faktorer som tilbakemeldinger, mestringsfølelse og sosial støtte \parencite[s.~152]{einarsen2023}. Et kompromiss som forener organisasjonens koordineringsbehov med medarbeidernes behov for autonomi og forutsigbarhet, vil være avgjørende for å gjenopprette den psykologiske kontrakten og sikre både motivasjon og fremtidig måloppnåelse.
